\documentclass[12pt, a4paper]{ctexart}
\usepackage[margin=2cm]{geometry}
\usepackage{enumitem}
\usepackage{fancyhdr}
\usepackage{amsmath, amssymb}
\usepackage{graphicx}
\usepackage{xcolor}

\setCJKmainfont{Hiragino Sans GB}
\setCJKsansfont{Heiti SC}

\pagestyle{fancy}
\fancyhf{}
\fancyhead[L]{八年级上册数学测试卷}
\fancyhead[R]{第~\thepage~页}
\fancyfoot[C]{}
\renewcommand{\headrulewidth}{0.4pt}

\begin{document}

\begin{center}
    {\LARGE\bfseries 八年级上册数学测试卷}\\[0.5cm]
    {\large 考试时间：120分钟 \quad 满分：120分}\\[0.3cm]
    {\normalsize 注意：请将答案写在答题卡上，写在试卷上无效。}
\end{center}

\vspace{0.5cm}

{\large\bfseries 一、选择题（本大题共10小题，每小题3分，共30分。每小题只有一个正确选项。）}

\begin{enumerate}[label=\arabic*.]

\item 已知三角形的一个角为$50^{\circ}$，另两个角分别为$2x+10$和$2x-10$（单位：度），则$x$的值为（\hspace{1cm}）
\begin{enumerate}[label=\Alph*.]
\quad \textbf{A.}~$40$ \quad \textbf{B.}~$45$ \quad \textbf{C.}~$50$ \quad \textbf{D.}~$55$
\end{enumerate}

\item 下列各组线段中，能构成三角形的是（\hspace{1cm}）
\begin{enumerate}[label=\Alph*.]
\quad \textbf{A.}~$3,4,5$ \quad \textbf{B.}~$1,2,3$ \quad \textbf{C.}~$2,3,6$ \quad \textbf{D.}~$4,5,10$
\end{enumerate}

\item 下列图形中，是轴对称图形的是（\hspace{1cm}）
\begin{enumerate}[label=\Alph*.]
\quad \textbf{A.}~平行四边形 \quad \textbf{B.}~等腰梯形 \quad \textbf{C.}~一般三角形 \quad \textbf{D.}~一般梯形
\end{enumerate}

\item 计算$(x+2)(x-3)$的结果是（\hspace{1cm}）
\begin{enumerate}[label=\Alph*.]
\quad \textbf{A.}~$x^{2}-x-6$ \quad \textbf{B.}~$x^{2}+x-6$ \quad \textbf{C.}~$x^{2}-x+6$ \quad \textbf{D.}~$x^{2}+x+6$
\end{enumerate}

\item 下列多项式中，不能用平方差公式分解因式的是（\hspace{1cm}）
\begin{enumerate}[label=\Alph*.]
\quad \textbf{A.}~$x^{2}-1$ \quad \textbf{B.}~$4a^{2}-9b^{2}$ \quad \textbf{C.}~$-x^{2}+y^{2}$ \quad \textbf{D.}~$x^{2}+4$
\end{enumerate}

\item 化简$\dfrac{x^{2}-4}{x^{2}-4x+4}$的结果是（\hspace{1cm}）
\begin{enumerate}[label=\Alph*.]
\quad \textbf{A.}~$\dfrac{x+2}{x-2}$ \quad \textbf{B.}~$\dfrac{x-2}{x+2}$ \quad \textbf{C.}~$1$ \quad \textbf{D.}~$-1$
\end{enumerate}

\item 等腰三角形的两边长分别为$5$和$10$，则该三角形的周长为（\hspace{1cm}）
\begin{enumerate}[label=\Alph*.]
\quad \textbf{A.}~$20$ \quad \textbf{B.}~$25$ \quad \textbf{C.}~$20$或$25$ \quad \textbf{D.}~$15$
\end{enumerate}

\item 计算$(-2x^{2}y)^{3}$的结果是（\hspace{1cm}）
\begin{enumerate}[label=\Alph*.]
\quad \textbf{A.}~$-6x^{6}y^{3}$ \quad \textbf{B.}~$-8x^{6}y^{3}$ \quad \textbf{C.}~$8x^{6}y^{3}$ \quad \textbf{D.}~$-8x^{5}y^{3}$
\end{enumerate}

\item 如图，在$\triangle ABC$中，$AB=AC$，$D$是$BC$的中点，$\angle BAD=35^{\circ}$，则$\angle C$的度数为（\hspace{1cm}）
\begin{enumerate}[label=\Alph*.]
\quad \textbf{A.}~$35^{\circ}$ \quad \textbf{B.}~$55^{\circ}$ \quad \textbf{C.}~$65^{\circ}$ \quad \textbf{D.}~$70^{\circ}$
\end{enumerate}

\item 已知等腰三角形的一个角为$80^{\circ}$，则它的另外两个角分别是（\hspace{1cm}）
\begin{enumerate}[label=\Alph*.]
\quad \textbf{A.}~$20^{\circ},80^{\circ}$ \quad \textbf{B.}~$50^{\circ},50^{\circ}$ \quad \textbf{C.}~$80^{\circ},20^{\circ}$ \quad \textbf{D.}~$50^{\circ},50^{\circ}$或$80^{\circ},20^{\circ}$
\end{enumerate}

\end{enumerate}

\vspace{0.5cm}

{\large\bfseries 二、填空题（本大题共6小题，每小题3分，共18分。）}

\begin{enumerate}[label=\arabic*., start=11]

\item 计算：$(a^{2}b)^{3}=$ \underline{\hspace{3cm}} 。

\item 三角形的两条边长分别为$3$和$7$，则第三边的取值范围是 \underline{\hspace{3cm}} 。

\item 分解因式：$x^{2}-5x+6=$ \underline{\hspace{3cm}} 。

\item 若等腰三角形的底边长为$6$，周长为$16$，则腰长为 \underline{\hspace{3cm}} 。

\item 计算：$\dfrac{x}{x-1}-\dfrac{1}{x-1}=$ \underline{\hspace{3cm}} 。

\item 如图，在$\triangle ABC$中，$\angle C=90^{\circ}$，$AD$是$\angle BAC$的平分线，$DE\perp AB$于$E$，若$AC=6$，$BC=8$，$AB=10$，则$DE=$ \underline{\hspace{3cm}} 。

\end{enumerate}

\vspace{0.5cm}

{\large\bfseries 三、解答题（本大题共8小题，共72分。解答应写出文字说明、演算步骤或证明过程。）}

\begin{enumerate}[label=\arabic*., start=17]

\item （8分）已知$\triangle ABC$中，$D$是$BC$的中点，$DE\perp AB$于$E$，$DF\perp AC$于$F$，且$DE=DF$。

求证：$\triangle BDE \cong \triangle CDF$。

\vspace{6cm}

\item （8分）先化简，再求值：$(2x+1)(2x-1)-4x(x-1)$，其中$x=\dfrac{1}{2}$。

\vspace{6cm}

\item （8分）化简求值：$\dfrac{x}{x^{2}-4}-\dfrac{1}{x^{2}-4x+4}$，其中$x=3$。

\vspace{6cm}

\item （10分）如图，在$\triangle ABC$中，$\angle C=90^{\circ}$，$AC=6$，$BC=8$，$CD\perp AB$于$D$。

（1）求$AB$的长；

（2）求$CD$的长。

\vspace{6cm}

\item （10分）如图，已知$AB=CD$，$\angle B=\angle C$，$BF=CE$。

求证：$\triangle ABF \cong \triangle DCE$。

\vspace{6cm}

\item （8分）已知$x^{2}-2x+m=(x-3)(x+n)$，其中$m$、$n$为常数，求$m+n$的值。

\vspace{6cm}

\item （10分）如图，在$\triangle ABC$中，$AB=AC$，$D$是$AB$的中点，$DE\parallel BC$交$AC$于$E$。

（1）求证：$DE$平分$AE$；

（2）若$AB=10$，$BC=6$，求$\triangle ADE$的周长。

\vspace{6cm}

\item （10分）某文具店笔记本每本$5$元，钢笔每支$y$元。小明买了$4$本笔记本和$6$支钢笔共花费$70$元。

（1）求每支钢笔的价格；

（2）若小明买$m$本笔记本和$n$支钢笔恰好花费$70$元，求$m$和$n$的关系式。

\vspace{6cm}

\end{enumerate}

\end{document}
